\documentclass[9pt,a4paper,twocolumn]{ctexart}

% --- Packages ---
\usepackage[top=1.2cm, bottom=1.2cm, left=1.2cm, right=1.2cm, columnsep=0.8cm]{geometry}
\usepackage{hyperref}
\usepackage{graphicx}
\usepackage{float}
\usepackage{longtable}
\usepackage{tabularx}
\usepackage{booktabs}
\usepackage{enumitem}
\usepackage{xcolor}
\usepackage{fancyhdr}
\usepackage{listings}
\usepackage{fontspec}
\usepackage{titlesec}
\usepackage{caption}

% --- Code block styling ---
\definecolor{codebg}{RGB}{245,245,245}
\definecolor{codeframe}{RGB}{220,220,220}
\definecolor{codekw}{RGB}{0,0,192}
\definecolor{codecomment}{RGB}{0,128,0}
\definecolor{codestring}{RGB}{163,21,21}
\definecolor{codenum}{RGB}{128,0,128}
\definecolor{codepunct}{RGB}{90,90,90}
\definecolor{badgegreen}{RGB}{34,139,34}
\definecolor{badgeblue}{RGB}{30,144,255}
\definecolor{badgeblack}{RGB}{50,50,50}

% --- Difficulty badges (rounded corners via tikz, pre-rendered in saveboxes) ---
\usepackage{tikz}
\newsavebox{\badgechubox}
\savebox{\badgechubox}{\tikz[baseline]{\node[fill=badgegreen,text=white,rounded corners=2.5pt,inner xsep=4pt,inner ysep=1.5pt,font=\sffamily\footnotesize\bfseries]{初级};}}
\newsavebox{\badgezhongbox}
\savebox{\badgezhongbox}{\tikz[baseline]{\node[fill=badgeblue,text=white,rounded corners=2.5pt,inner xsep=4pt,inner ysep=1.5pt,font=\sffamily\footnotesize\bfseries]{中级};}}
\newsavebox{\badgegaobox}
\savebox{\badgegaobox}{\tikz[baseline]{\node[fill=badgeblack,text=white,rounded corners=2.5pt,inner xsep=4pt,inner ysep=1.5pt,font=\sffamily\footnotesize\bfseries]{高级};}}
\newcommand{\lvchu}{\usebox{\badgechubox}}
\newcommand{\lvzhong}{\usebox{\badgezhongbox}}
\newcommand{\lvgao}{\usebox{\badgegaobox}}
\pdfstringdefDisableCommands{%
  \def\lvchu{[初级]}%
  \def\lvzhong{[中级]}%
  \def\lvgao{[高级]}%
}

\lstset{
  basicstyle=\ttfamily\scriptsize,
  keywordstyle=\color{codekw}\bfseries,
  commentstyle=\color{codecomment}\itshape,
  stringstyle=\color{codestring},
  backgroundcolor=\color{codebg},
  frame=single,
  rulecolor=\color{codeframe},
  breaklines=true,
  breakatwhitespace=false,
  breakindent=0pt,
  postbreak=\mbox{\textcolor{red}{$\hookrightarrow$}\space},
  numbers=left,
  numberstyle=\tiny\color{gray},
  columns=flexible,
  keepspaces=true,
  showstringspaces=false,
  tabsize=2,
  aboveskip=4pt,
  belowskip=4pt,
  extendedchars=true,
  literate={，}{{，}}1 {；}{{；}}1 {！}{{！}}1 {？}{{？}}1 {“}{{“}}1 {”}{{”}}1 {（}{{（}}1 {）}{{）}}1 {《}{{《}}1 {》}{{》}}1,
}

% --- Extra language definitions (no built-in listings support) ---
\lstdefinelanguage{json}{
  morekeywords={true,false,null},
  sensitive=true,
  morestring=[b]",
  comment=[l]{//},
  morecomment=[s]{/*}{*/},
}
\lstdefinelanguage{yaml}{
  morekeywords={true,false,null,yes,no,on,off},
  sensitive=false,
  morecomment=[l]{\#},
  morestring=[b]",
  morestring=[d]',
}
\lstdefinelanguage{http}{
  morekeywords={GET,POST,PUT,DELETE,PATCH,HEAD,OPTIONS,HTTP,Host,Accept,Authorization,Content-Type,Content-Length,User-Agent},
  sensitive=true,
  morecomment=[l]{\#},
  morestring=[b]",
}
\lstdefinelanguage{TypeScript}{
  morekeywords={abstract,any,as,async,await,break,case,catch,class,const,continue,debugger,declare,default,delete,do,else,enum,export,extends,false,finally,for,from,function,get,if,implements,import,in,instanceof,interface,keyof,let,module,namespace,never,new,null,number,of,package,private,protected,public,readonly,return,set,static,string,super,switch,symbol,this,throw,true,try,type,typeof,undefined,unknown,var,void,while,with,yield,Record,Array,Promise,AsyncIterable,ReadableStream,Date,Error,Object,JSON},
  sensitive=true,
  morecomment=[l]{//},
  morecomment=[s]{/*}{*/},
  morestring=[b]",
  morestring=[b]',
  morestring=[b]`{`},
}

% --- Compact spacing ---
\linespread{0.95}
\setlength{\parskip}{1pt plus 1pt minus 1pt}
\setlength{\parindent}{0pt}

% --- Table styling ---
\renewcommand{\arraystretch}{1.1}

% --- Heading styling (numbered) ---
\titleformat{\section}{\normalsize\bfseries}{\thesection\quad}{0em}{}
\titlespacing{\section}{0pt}{6pt}{2pt}
\titleformat{\subsection}{\small\bfseries}{\thesubsection\quad}{0em}{}
\titlespacing{\subsection}{0pt}{4pt}{1pt}
\titleformat{\subsubsection}{\small\bfseries}{\thesubsubsection\quad}{0em}{}
\titlespacing{\subsubsection}{0pt}{3pt}{1pt}

% --- Page style ---
\pagestyle{fancy}
\fancyhf{}
\fancyhead[L]{\small\emph{\leftmark}}
\fancyhead[R]{\small\thepage}
\renewcommand{\headrulewidth}{0.4pt}
\renewcommand{\sectionmark}[1]{}  % prevent sections from overwriting leftmark

% --- Hyperref setup ---
\hypersetup{
  colorlinks=true,
  linkcolor=blue,
  urlcolor=blue,
  citecolor=blue,
}

% --- Caption format ---
\DeclareCaptionFormat{myformat}{\small\bfseries #1#2#3}
\captionsetup{format=myformat}

\begin{document}

\title{上下文工程：预算、裁剪、缓存与组装}
\date{}
\maketitle
\markboth{Java 版 Agent 知识}{ Java 版 Agent 知识}

\begin{quote}
语言策略：\texttt{code}。正文三语言一致；上下文组装器按 Java 生态实现。 地图节点：\texttt{04 开发与工程化 / 上下文工程}。配套文档：\href{03-Harness工程.md}{Harness Engineering}、\href{../02-Agent/02-Agent记忆系统.md}{Agent 记忆系统}。
\end{quote}



\section*{0. 结论先说}

\begin{enumerate}[itemsep=1pt, topsep=2pt]
  \item Agent 效果瓶颈常在 Context Engineering，不在模型。
  \item 每一轮喂给模型的上下文应当是\textbf{预算内、按优先级排序、边界清晰}的组装结果。
  \item 上下文工程 = 预算 + 选择 + 压缩 + 组装 + 缓存 + 版本化。
\end{enumerate}


\section*{1. 上下文预算公式}


\begin{lstlisting}
可用上下文 = 模型窗口 - 保留输出 Token - System Prompt - 工具 Schema
每轮动态预算 = 可用上下文 - 历史对话预算 - 检索预算 - 记忆预算 - 安全余量
\end{lstlisting}



规则：保留输出 Token 按任务最大回复长度计算；工具 Schema 只注入本轮白名单工具；安全余量默认留 10\%。



\section*{2. 内容优先级}

{\footnotesize
\begin{tabular}{|p{\dimexpr 0.083\columnwidth-2\tabcolsep\relax}|p{\dimexpr 0.389\columnwidth-2\tabcolsep\relax}|p{\dimexpr 0.528\columnwidth-2\tabcolsep\relax}|}
\hline
\textbf{优先级} & \textbf{内容} & \textbf{处理策略} \\
\hline
P0 & 任务目标、验收标准、安全规则 & 永不被裁剪 \\
\hline
P1 & 最近 N 轮观察 & 保留最近，更早做摘要 \\
\hline
P2 & 工具返回 & 截断长文本，保留结构化字段 \\
\hline
P3 & 检索片段 & 只保留 rerank 后的 top-k \\
\hline
P4 & 历史记忆 & 按相关度注入，低相关不注入 \\
\hline
\end{tabular}
}



\section*{3. 压缩与裁剪}

\begin{itemize}[itemsep=1pt, topsep=2pt]
  \item 工具返回先提取结构化摘要，再决定是否保留原文；
  \item 历史对话超过预算时，把早期轮次压缩成状态摘要；
  \item 代码和日志用行号范围截断，不从中部硬切；
  \item 外部内容统一包裹为数据区，防止指令注入；
  \item 达到 40\% 上下文利用率后，优先裁剪而不是继续塞入。
\end{itemize}


\section*{4. 缓存与版本化}

\begin{itemize}[itemsep=1pt, topsep=2pt]
  \item System Prompt 固定部分放前缀，命中 Prompt Cache；
  \item 工具 Schema 按 toolset\_version 缓存；
  \item 检索缓存键 = query 归一化 + 知识库版本 + 租户；
  \item 每次上下文模板变更要 bump \texttt{prompt\_version}，并跑评测回归。
\end{itemize}


\section*{5. 组装顺序}


\begin{lstlisting}
System Prompt（版本化）
→ 任务卡与验收标准
→ 本轮白名单工具 Schema
→ 相关记忆（低相关不注入）
→ 检索片段（标注来源）
→ 最近观察（按预算截断）
→ 用户当前输入（隔离包裹）
\end{lstlisting}



\section*{6. Java 版上下文组装器}


\begin{lstlisting}[language=Java]
public record ContextBudget(int windowTokens, int reservedOutput, int systemPrompt,
                            int toolSchemas, int safetyReserve) {
    public int available() {
        return windowTokens - reservedOutput - systemPrompt - toolSchemas - safetyReserve;
    }
}

public record ContextSection(String name, int priority, int tokens, String content) {}

public class ContextAssembler {
    public static String assemble(ContextBudget budget, List<ContextSection> sections, int maxTokens) {
        int cap = Math.min(maxTokens, budget.available());
        StringBuilder sb = new StringBuilder();
        int used = 0;
        for (ContextSection s : sections.stream()
                .sorted(Comparator.comparingInt(ContextSection::priority)).toList()) {
            if (used + s.tokens() > cap && s.priority() > 1) continue;
            sb.append("<!-- section:").append(s.name()).append(" -->\n")
              .append(s.content()).append("\n");
            used += s.tokens();
        }
        return sb.toString();
    }
}
\end{lstlisting}



\section*{7. 上线检查单}

\begin{itemize}[itemsep=1pt, topsep=2pt]
  \item [ ] 每次调用都计算预算，不超模型窗口；
  \item [ ] 工具 Schema 只注入白名单工具；
  \item [ ] 长文本按优先级裁剪，不丢失任务目标；
  \item [ ] 外部数据有数据边界标记；
  \item [ ] Prompt 版本与评测结果绑定；
  \item [ ] 缓存命中率、单任务 Token、上下文利用率有监控。
\end{itemize}


\section*{8. 延伸阅读}

\begin{itemize}[itemsep=1pt, topsep=2pt]
  \item \href{../02-Agent/02-Agent记忆系统.md}{Agent 记忆系统}
  \item \href{04-大模型网关.md}{大模型网关}
  \item \href{../05-评测与质量/01-Agent评测体系.md}{Agent 评测体系}
\end{itemize}

\end{document}
